\id{IRSTI 06.58.30}{}

\begin{header}
\swa{G.P.~Koptayeva, A.B.~Mukhamedkhanova, A.N.~Issakhmetova, G.A.~Myrzalieva, A.A.~Mutaliyeva, D.N.~Kelesbayev}{THE NATURAL RESOURCE CURSE IN KAZAKHSTAN: ECONOMIC AND STRUCTURAL PROBLEMS AND SOLUTION PROPOSALS}

\tsp{1}G.P.~Koptayeva,
\tsp{2}A.B.~Mukhamedkhanova\envelope,
\tsp{3}A.N.~Issakhmetova,
\tsp{1}G.A.~Myrzalieva,
\tsp{4}A.A.~Mutaliyeva,
\tsp{5}D.N.~Kelesbayev
\end{header}

\begin{affil}
\tsp{1}Miras University, Shymkent, Kazakhstan,

\tsp{2}M. Auezov South Kazakhstan University, Shymkent, Kazakhstan,

\tsp{3}Central-Asian Innovation University, Shymkent, Kazakhstan,

\tsp{4}Regional Innovation University, Shymkent, Kazakhstan,

\tsp{5}Akhmet Yassawi University, Turkestan, Kazakhstan

\corrauthor{Correspondent-author: Dia-2808@mail.ru}
\end{affil}

This study aims to analyze the structural vulnerabilities caused by
Kazakhstan' s natural resource-based economic structure
and to identify transformation requirements for sustainable development.
The research conducted a literature review based on secondary data and
utilized SWOT and PEST analyses to assess the country' s
internal and external dynamics. The findings demonstrate that
Kazakhstan' s economic growth is largely dependent on oil
and natural gas revenues, leading to fundamental problems such as a lack
of economic diversification, inadequate institutional capacity, and high
vulnerability to external shocks. Theoretically, consistent with the
natural resource curse, the Dutch disease, and the linkage theory,
Kazakhstan has achieved limited success in transforming its natural
resource wealth into sustainable development. Comparisons with the
examples of Norway, Canada, and Botswana reveal that successful
countries have transformed resource dependency into an opportunity
through strong institutional structures, transparent revenue management,
and economic diversification. For Kazakhstan, institutional reforms,
human capital investments, and green economy strategies are highlighted
as priority policy areas. Based on the data provided by SWOT and PEST
analyses, the research systematically examines the main threats and
potential opportunities Kazakhstan faces on its path to sustainable
development. The results point to the need for multidimensional policy
interventions to transform natural resource wealth into a development
advantage. In this context, recommendations are presented as strategic
roadmaps for policymakers and academia.

{\bfseries Keywords:} natural resource curse, economic diversification,
institutional capacity, SWOT analysis, PEST analysis, sustainable
development, Dutch disease, Kazakhstan economy.

\begin{header}
ҚАЗАҚСТАНДАҒЫ ТАБИҒИ РЕСУРСТАРДЫҢ ҚАРҒЫСЫ: ЭКОНОМИКАЛЫҚ ЖӘНЕ ҚҰРЫЛЫМДЫҚ МӘСЕЛЕЛЕР ЖӘНЕ ШЕШІМ ҰСЫНЫСТАРЫ

\tsp{1}Г.П.~Коптаева,
\tsp{2}А.Б.~Мухамедханова\envelope,
\tsp{3}А.Н.~Исахметова,
\tsp{1}Г.А.~Мырзалиева,
\tsp{4}А.А.~Муталиева,
\tsp{5}Д.Н.~Келесбаев
\end{header}

\begin{affil}
\tsp{1}Мирас университеті, Шымкент, Қазақстан,

\tsp{2}М.Әуезов атындағы Оңтүстік Қазақстан университеті, Шымкент, Қазақстан,

\tsp{3}Орталық Азия инновациялық университеті, Шымкент, Қазақстан,

\tsp{4}Аймақтық инновациялық университеті, Шымкент, Қазақстан,

\tsp{5}Ахмет Ясауи университеті, Түркістан, Қазақстан,

e-mail: Dia-2808@mail.ru
\end{affil}

Бұл зерттеудің мақсаты-Қазақстанның табиғи ресурстарға негізделген
экономикалық құрылымынан туындаған құрылымдық осалдықтарды талдау және
тұрақты даму үшін трансформация талаптарын анықтау. Зерттеу барысында
екінші реттік деректерге негізделген әдебиеттерге шолу жасалды және
елдің ішкі және сыртқы динамикасын бағалау үшін SWOT және PEST
талдаулары қолданылды. Зерттеу нәтижелері Қазақстанның экономикалық өсуі
көбінесе мұнай мен табиғи газ кірістеріне тәуелді екенін көрсетеді, бұл
экономикалық әртараптандырудың болмауы, институционалдық әлеуеттің
жеткіліксіздігі және сыртқы соққыларға жоғары осалдығы сияқты іргелі
мәселелерге әкеледі. Теориялық тұрғыдан алғанда, табиғи ресурстардың
қарғысына, голланд ауруына және байланыс теориясына сәйкес, Қазақстан
өзінің табиғи ресурстарға байлығын тұрақты дамуға айналдыруда шектеулі
табысқа жетті. Норвегия, Канада және Ботсвананың мысалдарымен салыстыру
табысты елдердің күшті институционалдық құрылымдар, ашық кірістерді
басқару және экономикалық әртараптандыру арқылы ресурстарға тәуелділікті
мүмкіндікке айналдырғанын көрсетеді. Қазақстан үшін институционалдық
реформалар, адами капиталға инвестициялар және жасыл экономика
стратегиялары басым саясат салалары ретінде белгіленген. SWOT және PEST
талдаулары ұсынған деректерге сүйене отырып, зерттеу Қазақстанның
тұрақты даму жолында кездесетін негізгі қауіптер мен әлеуетті
мүмкіндіктерді жүйелі түрде зерттейді. Нәтижелер табиғи ресурстар
байлығын даму артықшылығына айналдыру үшін көп өлшемді саяси
араласулардың қажеттілігін көрсетеді. Осыған байланысты ұсыныстар
саясаткерлер мен академиялық топтар үшін стратегиялық жол карталары
ретінде ұсынылған.

{\bfseries Түйін сөздер:} табиғи ресурстардың қарғысы, экономикалық
әртараптандыру, институционалдық әлеует, SWOT талдауы, PEST талдауы,
тұрақты даму, голланд ауруы, Қазақстан экономикасы.

\begin{header}
ПРОКЛЯТИЕ ПРИРОДНЫХ РЕСУРСОВ В КАЗАХСТАНЕ: ЭКОНОМИЧЕСКИЕ И СТРУКТУРНЫЕ ПРОБЛЕМЫ И ПРЕДЛОЖЕНИЯ ПО РЕШЕНИЮ

\tsp{1}Г.П.~Коптаева,
\tsp{2}А.Б.~Мухамедханова,
\tsp{3}А.Н.~Исахметова,
\tsp{1}Г.А.~Мырзалиева,
\tsp{4}А.А.~Муталиева,
\tsp{5}Д.Н.~Келесбаев\envelope
\end{header}

\begin{affil}
\tsp{1}Университет Мирас, Шымкент, Казахстан,

\tsp{2}Южно-Казахстанский университет имени М.Ауэзова, Шымкент, Казахстан,

\tsp{3}Центрально-Азиатский инновационный университет, Шымкент, Казахстан,

\tsp{4}Региональный инновационный университет, Шымкент, Казахстан,

\tsp{5}Университет Ахмеда Ясави, Туркестан, Казахстан,

e-mail: Dia-2808@mail.ru
\end{affil}

Целью данного исследования является анализ структурной уязвимости,
обусловленной ресурсоориентированной структурой экономики Казахстана, и
определение необходимых преобразований для обеспечения устойчивого
развития. В рамках исследования был проведен обзор литературы на основе
вторичных данных, а также применены SWOT- и PEST-анализы для оценки
внутренней и внешней динамики страны. Результаты показывают, что
экономический рост Казахстана в значительной степени зависит от доходов
от нефти и природного газа, что приводит к таким фундаментальным
проблемам, как отсутствие экономической диверсификации, недостаточный
институциональный потенциал и высокая уязвимость к внешним шокам.
Теоретически, в соответствии с «ресурсным проклятием», «голландской
болезнью» и теорией взаимосвязей, Казахстан добился ограниченных успехов
в преобразовании своего богатства от природных ресурсов в устойчивое
развитие. Сравнение с примерами Норвегии, Канады и Ботсваны показывает,
что успешные страны преобразовали ресурсную зависимость в возможности
благодаря сильным институциональным структурам, прозрачному управлению
доходами и экономической диверсификации. Для Казахстана
институциональные реформы, инвестиции в человеческий капитал и стратегии
зеленой экономики выделены в качестве приоритетных направлений политики.
На основе данных SWOT- и PEST-анализов исследование систематически
рассматривает основные угрозы и потенциальные возможности, с которыми
сталкивается Казахстан на пути к устойчивому развитию. Результаты
указывают на необходимость многомерных политических мер для
преобразования богатства природных ресурсов в преимущество в области
развития. В этом контексте рекомендации представлены в виде
стратегических дорожных карт для политиков и академических кругов.

{\bfseries Ключевые слова:} проклятие природных ресурсов, экономическая
диверсификация, институциональный потенциал, SWOT-анализ, PEST-анализ,
устойчивое развитие, голландская болезнь, экономика Казахстана.

\begin{multicols}{2}
{\bfseries Introduction.} Natural resources have historically played an
important role in the economic development of many countries. However,
in some cases, resource abundance may hinder long-term development by
creating economic distortions, institutional weaknesses, and
environmental challenges. This phenomenon is widely known in the
literature as the ``resource curse'' {[}1,2{]}.

Kazakhstan is one of the most resource-rich countries in Central Asia.
Although natural resources have been an important source of economic
growth, excessive dependence on the extractive sector has generated
structural economic, social, and environmental challenges {[}3, 4{]}.

After gaining independence, Kazakhstan faced major structural challenges
during its transition to a market economy. Resource revenues became a
key driver of economic growth; however, this dependence also contributed
to structural imbalances, limited diversification, and institutional
challenges {[}5- 7{]}.

The aim of this study is to examine Kazakhstan' s current
economic structure and propose concrete strategic steps to reduce its
dependence on natural resources and establish a sustainable economic
structure. Within this framework, the study emphasizes that areas such
as industrialization, education reforms, green transformation, foreign
trade strategies, and institutional development offer potential
solutions for Kazakhstan. Successful implementation of these strategies
can contribute to Kazakhstan' s long-term development
goals by increasing its economic diversification.

To this end, strategic analytical approaches are proposed to support
Kazakhstan in adopting a more balanced and sustainable development model
by reducing its dependence on natural resources. In this context, SWOT
(Strengths, Weaknesses, Opportunities, Threats) and PEST (Political,
Economic, Social, Technological) analytical frameworks are employed in
this study.

Within this framework, the study highlights several key policy areas
including economic diversification, human capital development, green
transformation, and international economic integration.

The following section outlines the theoretical foundations of the study.
The relationship between natural resources and economic performance has
been the subject of theoretical and empirical investigation for many
years. Until the mid-20th century, the prevailing view was that abundant
natural resources would foster development by accelerating
industrialization. However, the Singer--Prebisch thesis a significant
critique of this optimistic view. According to the Singer-Prebisch
thesis, countries relying on primary commodity exports gradually become
dependent on industrialized core countries due to unfavorable trends in
their terms of trade. In other words, peripheral countries exporting raw
materials find themselves increasingly disadvantaged by core countries
importing processed industrial goods and face structural obstacles to
development. These early discussions first highlighted the risks of
natural resource-based growth strategies.

Modern literature has developed a number of theoretical approaches to
explain the mechanisms by which natural resource shocks affect the
economy. One of these approaches, the Dutch disease theory, suggests
that large inflows of foreign currency from natural resource exports may
appreciate the domestic currency and reduce the competitiveness of other
sectors. This appreciation of the local currency leads to higher export
prices for tradable sectors such as manufacturing and agriculture,
leading to a contraction in production. Another approach, linkage
theory, draws on the work of Albert Hirschman and emphasizes that
industrialization will be limited if the natural resource sector cannot
establish forward and backward linkages with other sectors {[}8{]}.
Strong linkages imply that natural resource revenues are channeled into
investments in sectors such as manufacturing and services, thereby
fostering economic expansion. If the natural resource sector remains an
enclave within the economy and does not support other sectors, its
growth potential is limited. Indeed, as some countries like Norway have
successfully demonstrated, integrating natural resource sectors like oil
with other sectors within a country can mitigate perceived negative
effects and create a balanced development model.

Empirical findings on the natural resource curse are quite
controversial. Many studies have found that natural resource abundance
is negatively correlated with economic growth {[}1,4{]}. However, some
argue in the literature that this relationship is not inevitable. Some
economists argue that resource abundance can be a boon, that is, it can
support growth, when the right policies are implemented {[}9,10{]}.
Conversely, many other studies point to resource abundance as the
primary cause of economic stagnation or contraction. Some researchers
have even argued that the natural resource curse is the exception, not
the rule: For example, Alssadek and Benhin {[}11{]} note that not all
resource-rich countries experience the curse, while Wang et al. {[}12{]}
and Herb {[}13{]} find no statistically significant link between natural
resource abundance and low growth. However, the common thread across all
these debates is the consensus on the decisive role of institutional
factors. As nearly all researchers agree, strong and transparent
institutions and a stable, diversified economic structure can mitigate
the potential negative impacts of natural resources. Recent studies have
shown that some countries (like Norway) have been able to leverage their
wealth of natural resources thanks to favorable institutional
environments, while others (like Nigeria) face economic hardship despite
their abundance of resources due to weak institutions. This has
contributed to a more balanced perspective in the literature.

Although numerous studies examine the macroeconomic effects of resource
dependence, fewer studies integrate strategic analytical tools to
evaluate both internal and external drivers of structural
transformation.

The scientific novelty of this study lies in the integration of SWOT and
PEST analytical frameworks with macroeconomic trend analysis in order to
assess Kazakhstan's resource dependency and identify structural
transformation pathways.

This research contributes to the existing literature in three ways:

First, it provides an updated structural analysis of Kazakhstan's
resource-based economy using macroeconomic data covering the period from
1992 to 2024.

Second, the study integrates strategic management tools (SWOT and PEST
analysis) with economic development theory to diagnose structural
weaknesses and policy priorities.

Third, it develops a policy-oriented analytical framework that
identifies key transformation pathways for Kazakhstan's transition from
a resource-dependent economy to a diversified and sustainable
development model.

{\bfseries Materials and methods.} After gaining independence in 1991,
Kazakhstan experienced the throes of transitioning to a market economy
until 1998. While real gross domestic product increased significantly
with the rise in oil prices from 2000 onward, a structure dependent on
natural resources emerged in terms of production, employment, and
foreign trade. Therefore, this section will present the macroeconomic
indicators of the Kazakh economy and analyze its dependence on natural
resources in terms of production, employment, and foreign trade. The
purpose of this analysis is to reveal the impact of production structure
on the factors leading to the natural resource curse. In this context, a
combination of macroeconomic indicators and Brent oil prices will
provide a more useful analysis of the Kazakh economy (see Table 1).
\end{multicols}

\tcap{Table 1 - Some Macroeconomic Indicators of Kazakhstan}
\begin{longtblr}[
  label = none,
  entry = none,
  note{} = {\emph{Source: {[}14,15{]}}},
]{
  cells = {c},
  cells = {font = \small},
  row{1} = {font=\bfseries\small},
  hlines,
  vlines,
}
Years & GDP (Billion Dollars) & Growth Rate (\%) & Inflation (\%) & Unemployment (\%) \\
1992  & 3                     & -2,9             & 1381           & n/a               \\
1993  & 5                     & -9,2             & 1662,3         & 0,5               \\
1994  & 12                    & -12,6            & 854,6          & 7,5               \\
1995  & 17                    & -8,2             & 60,4           & 11,0              \\
1996  & 21                    & 0,5              & 28,7           & 13,0              \\
1997  & 22                    & 1,7              & 11,2           & 13,0              \\
1998  & 22                    & -1,9             & 1,9            & 13,1              \\
1999  & 17                    & 2,7              & 17,8           & 13,5              \\
2000  & 18                    & 9,8              & 9,8            & 12,8              \\
2001  & 22                    & 13,5             & 6,4            & 10,4              \\
2002  & 25                    & 9,8              & 6,6            & 9,3               \\
2003  & 31                    & 9,3              & 6,8            & 8,8               \\
2004  & 43                    & 9,6              & 6,7            & 8,4               \\
2005  & 57                    & 9,7              & 7,6            & 8,1               \\
2006  & 81                    & 10,7             & 8,4            & 7,8               \\
2007  & 105                   & 8,9              & 18,8           & 7,3               \\
2008  & 133                   & 3,3              & 9,5            & 6,6               \\
2009  & 115                   & 1,2              & 6,2            & 6,6               \\
2010  & 148                   & 7,3              & 7,8            & 5,8               \\
2011  & 193                   & 7,5              & 7,4            & 5,4               \\
2012  & 208                   & 5,0              & 6,0            & 5,3               \\
2013  & 237                   & 5,9              & 4,8            & 5,2               \\
2014  & 221                   & 4,3              & 7,4            & 5,1               \\
2015  & 184                   & 1,0              & 13,6           & 5,0               \\
2016  & 137                   & 0,9              & 8,5            & 4,9               \\
2017  & 167                   & 3,9              & 7,1            & 4,9               \\
2018  & 179                   & 4,1              & 5,3            & 4,8               \\
2019  & 182                   & 4,5              & 5,4            & 4,8               \\
2020  & 171                   & -2,6             & 7,5            & 4,9               \\
2021  & 197                   & 4,1              & 8,4            & 4,9               \\
2022  & 226                   & 3,3              & 20,3           & 4,9               \\
2023  & 263                   & 5,1              & 9,8            & 4,8               \\
2024  & 293                   & 3,5              & 8,0            & 4,8               
\end{longtblr}

\begin{multicols}{2}
As shown in Table 1, Kazakhstan's macroeconomic data
from 1992 to 2024 reveal the transformation and structural
vulnerabilities of natural resource-based economies. Having newly gained
independence in the early 1990s, the country experienced major economic
contractions during the transitional economy. Negative growth,
hyperinflation (e.g., 1662.3\% in 1993) and double-digit unemployment
rates are notable during the 1993-1995 period. Thanks to the growth
model based on oil and natural gas exports in the post-2000 period,
Kazakhstan' s GDP has increased significantly. GDP, which
was \$18 billion in 2000, rose to \$237 billion in 2013. During this
period, growth rates varied between 9-13\%; however, inflation also
remained in double-digits at times. The unemployment rate, on the other
hand, showed a steady downward trend. After 2014, the growth rate slowed
down due to the decline in oil prices, the ruble crisis and global
fluctuations; GDP declined to \$184 billion in 2015. Inflation rose to
13.6\% in 2015, and growth nearly stalled (1.0\%). This demonstrates
that the Kazakh economy remains highly vulnerable to external shocks.
After contracting by 2.6\% in 2020 due to the pandemic, the economy
recovered in subsequent years. Growth was recorded at 5.1\% in 2023,
with GDP reaching \$263 billion. However, inflation reached a remarkable
20.3\% in 2022, posing risks to the sustainability of price stability.
Overall, these data demonstrate that Kazakhstan' s
natural resource-based growth has provided short-term prosperity, but
the economic structure remains vulnerable to external shocks. While GDP
has increased, inflation has fluctuated, and weak institutional
structures and production diversification have exacerbated fragility.
Sustainable growth requires not only resource revenues but also
structural transformation in the industrial and service sectors.

Accordingly, production diversification aimed to make
Kazakhstan' s economic structure more balanced and
sustainable. Growth in the agricultural and manufacturing sectors has
helped to gradually reduce the share of natural resource revenues in
GDP. Thanks to these diversification efforts, Kazakhstan has partially
transitioned from an energy-based economy to a more diversified one.
Beginning in the 2010s, driven by production diversification strategies,
the share of Kazakhstan' s natural resource revenues in
GDP has gradually declined, while the contributions of the agricultural
and manufacturing sectors have increased. At the same time, the
country' s economy has become less dependent on external
resources, and the share of natural resource-based revenues has
decreased (see Figure 1).

According to Figure 1, Kazakhstan' s economic growth is
largely dependent on natural resources, particularly the oil and natural
gas sector. Over the period spanning 1992-2024, the share of natural
resource revenues in GDP has fluctuated depending on the
country' s economic policies, global market conditions,
and fluctuations in energy prices. Between these years, the Kazakh
economy has followed a natural resource-based growth model, with the oil
and natural gas sector serving as a key driver of the
country' s economic development. The share of natural
resource revenues in GDP peaked at over 66\% in the mid-2000s, but this
ratio has declined in subsequent years due to economic diversification
policies and global market fluctuations. Although a significant decline
was observed in the ratio of natural resource revenues to GDP in 2015,
this ratio has risen again since 2020. It can be said that as of 2024,
the economic contribution of natural resource revenues remains
significant, but efforts to transform Kazakhstan' s
economy into a more sustainable and diversified structure are ongoing.

After gaining independence in 1991, Kazakhstan implemented a series of
reforms and strategies to achieve macroeconomic stability. The
country' s primary objectives were to transition to a
market economy, support economic growth, and maintain financial
stability. Within this framework, the government developed various
macroeconomic policies to control inflation, ensure fiscal discipline,
and reduce budget deficits. Table 2 below lists the main macroeconomic
policies implemented by Kazakhstan to control inflation, ensure fiscal
discipline, and reduce budget deficits.
\end{multicols}

\fig[0.7\textwidth]{e/image14}[Fig.1 - GDP/Ratio of Natural Resource Revenues to
GDP\\\normalfont{\emph{Source: {[}14, 15{]}}}]

\tcap{Table 2 - Macroeconomic Policies Implemented by Kazakhstan}
\begin{longtblr}[
  label = none,
  entry = none,
  note{} = {\emph{Note - Compiled by the authors}},
]{
  width = \linewidth,
  colspec = {Q[227]Q[715]},
  cells = {font = \small},
  row{1} = {c,font=\bfseries\small},
  column{1} = {c},
  hlines,
  vlines,
}
Policies                                  & Explanations                                                                                                                                     \\
Monetary Policy Improvements              & To control inflation, the Central Bank increased interest rates and tightened the money supply.                                                  \\
Budget Discipline Practices               & The government implemented policies that strengthened budget discipline to balance expenditures and revenues.                                    \\
Fiscal Reforms                            & Reforms were implemented to review the tax system, broaden the tax base, and increase tax compliance.                                            \\
Inflation Targeting                       & The Central Bank set a specific inflation target and shaped its monetary policies to achieve this target.                                        \\
Restructuring Social Assistance           & To reduce budget deficits, social assistance programs were reviewed, and regulations were implemented to ensure more efficient use of resources. \\
Expenditure Cuts                          & Public spending was cut and redirected to priority projects.                                                                                     \\
Policies to Encourage Domestic Production & Various incentives were provided to encourage domestic production, reduce foreign exchange demand, and reduce import dependency.                 \\
Control over Exchange Rates               & The Central Bank intervened using its foreign exchange reserves to control exchange rate fluctuations.                                           \\
Use of International Financing Sources    & The aim was to secure and restructure financial resources through cooperation with international financial institutions.                         \\
Economic Stabilization Programs           & Various economic stability programs were developed to control inflation and ensure fiscal discipline.                                            
\end{longtblr}

\begin{multicols}{2}
According to Table 2, following its independence, Kazakhstan pursued
structural transformation through rapid economic reforms, transitioning
to a market economy through privatization and tight monetary and fiscal
policies. The Central Bank implemented policies to combat inflation and
stabilize the exchange rate, while fiscal management targeted budget
discipline through tax reforms and incentives. Oil and natural gas
revenues supported economic growth, but excessive dependence on these
resources highlighted the need for economic diversification. The country
aimed to diversify its economic structure by encouraging foreign direct
investment and strengthening its human capital through investments in
education, healthcare, and social security. Kazakhstan also implemented
incentive packages to maintain stability in the face of global crises
and pandemics and targeted sustainable growth with development plans for
2025. Based on this data and information, we can say that this research is
based entirely on secondary data sources. A comprehensive literature
review was conducted on the subject, examining academic sources (books,
peer-reviewed articles, and theses). Additionally, reports and datasets
published by international organizations (World Bank, IMF, UNDP, EBRD,
etc.) on the Kazakh economy were evaluated. Official statistics provided
by the National Statistical Institute of Kazakhstan, investor bulletins,
and economic analysis reports were also used in the data collection
process. This secondary data, compiled from diverse and reliable
sources, provided a solid basis for the study' s
analyses.

The conceptual framework of the research was established through a
systematic review of the relevant literature. Within this scope,
concepts such as the "natural resource curse" and "Dutch disease," which
explain the effects of natural resource wealth on economic development,
as well as theoretical approaches such as linkage theory, were examined
in depth. The literature review combined classic studies that shed light
on the subject with recent contemporary research. Thus, the theoretical
framework was structured in a comprehensive and balanced manner,
encompassing both traditional perspectives and new trends.

SWOT and PEST analyses were utilized in conjunction to assess the
current structure and transformation potential of the Kazakhstani
economy. SWOT analysis reveals a country' s internal
factors (strengths and weaknesses), while PEST analysis systematically
examines the political, economic, social, and technological elements
that constitute the external environment. These two methods were
integrated, addressing internal dynamics and external factors with a
holistic approach. The data used in the analyses were compiled from the
aforementioned empirical sources; key economic indicators, institutional
indices, and similar data were visualized through tables and graphs. The
findings were interpreted qualitatively in light of these tables and
graphs. Each SWOT component was linked to the relevant PEST elements,
providing a comprehensive analysis of Kazakhstan' s
strengths and weaknesses, as well as the opportunities and threats it
faces.

The study also adopted a comparative method to better assess
Kazakhstan' s economic structure, which is based on
natural resource revenue. Within this framework, Norway, Canada, and
Botswana, among the countries that stand out in managing their natural
resource wealth for development, were selected and examined. How these
countries manage their natural resource revenues, strengthen their
institutional structures, and their long-term development performance
are evaluated alongside Kazakhstan' s case. This
comparative analysis highlights the impact of different policy
approaches and institutional practices on economic development,
identifying positive examples and strategies that can serve as lessons
for Kazakhstan.

It' s important to note that this research has certain
limitations. First, because the study relies solely on secondary
sources, no primary data (field research, surveys, or interviews) was
included. Therefore, the findings are limited by the accuracy and scope
of existing datasets. Furthermore, the historical range and sectoral
coverage of the data used pose limitations that may impact the
generalizability of the analysis results. These limitations were
considered when interpreting the results, and the research findings were
carefully evaluated within this framework.

{\bfseries Results and discussion.} The results indicate that Kazakhstan's
economic structure exhibits several characteristics of the natural
resource curse. Excessive reliance on oil and gas revenues has limited
economic diversification and increased vulnerability to external shocks.
These findings are consistent with previous studies suggesting that
resource dependence may constrain long-term economic development.

The analysis also reveals signs of Dutch disease dynamics. Rapid growth
in the oil sector has weakened the competitiveness of other industries,
particularly manufacturing. In addition, weak linkages between the
extractive sector and other industries have limited broader industrial
development. These patterns also reflect structural dependency typical
of resource-exporting economies, where reliance on raw material exports
increases exposure to external economic fluctuations.

Although Kazakhstan's resource-based growth model has generated
significant economic expansion since the 2000s, it has also produced
structural imbalances. By 2023, more than 50\% of the country's exports
consisted of oil and gas, making economic performance highly sensitive
to global commodity prices. The dominance of the extractive sector has
also constrained the development of manufacturing and agriculture,
limiting the diversification of the economy.

From an institutional perspective, resource dependence has created
governance challenges, including weaker fiscal accountability and
limited incentives for institutional reforms. Furthermore, the
capital-intensive nature of the oil sector has restricted employment
generation and contributed to regional income disparities. Overall,
Kazakhstan's resource-based growth model has created economic,
institutional, and social imbalances that may hinder long-term
sustainable development.

The SWOT and PEST analyses further clarify the internal strengths and
weaknesses of Kazakhstan's economy, as well as the external
opportunities and threats affecting its development.

Kazakhstan's main strengths include its substantial natural resource
reserves, strategic geographic location between Europe and Asia, and
recent investments in infrastructure and human capital. These factors
provide a foundation for future economic transformation.

However, several weaknesses remain. The economy continues to depend
heavily on fossil fuel exports, while limited diversification and
institutional inefficiencies constrain sustainable growth. In addition,
governance challenges and uneven income distribution highlight
difficulties in converting resource wealth into broader social benefits.
\end{multicols}

\tcap{Table 3 - SWOT Analytical Synthesis Table}
\begin{longtblr}[
  label = none,
  entry = none,
  note{} ={\emph{Note - Compiled by the authors}},
]{
  width = \linewidth,
  colspec = {Q[213]Q[292]Q[437]},
  cells = {c},
  cells = {font = \small},
  row{1} = {font=\bfseries\small},
  hlines,
  vlines,
}
Analytical Dimension & Key Findings                     & Strategic Implication                              \\
Strengths            & Large natural resource reserves  & Can finance diversification if managed effectively \\
Weaknesses           & Lack of economic diversification & Structural reforms required                        \\
Opportunities        & Renewable energy potential       & Green transition strategy                          \\
Threats              & Oil price volatility             & Need for resilient economic structure              
\end{longtblr}

\begin{multicols}{2}
At the same time, several opportunities exist for economic
transformation. These include renewable energy development, digital
transformation, and deeper regional economic integration through
initiatives such as the Belt and Road Initiative. These factors could
support the diversification of Kazakhstan's economy and strengthen its
long-term competitiveness.

Nevertheless, Kazakhstan also faces significant external risks. These
include volatility in global commodity markets, climate-related
challenges, geopolitical uncertainties, and technological competition.
Such risks may affect the stability of export revenues and slow economic
modernization if appropriate policy measures are not implemented.

To summarize the strategic assessment, the key results of the SWOT
analysis are presented in Table 3. While Kazakhstan possesses
considerable natural resource wealth, the economy remains highly
dependent on the extractive sector. At the same time, renewable energy
development and technological modernization create opportunities for
diversification. However, global oil price volatility and external
economic shocks remain major risks. These findings highlight the
importance of diversification, institutional strengthening, and economic
resilience for achieving sustainable development.

The results of the SWOT and PEST analyses highlight the key challenges
and opportunities facing Kazakhstan in achieving sustainable
development. Based on these findings, several priority policy directions
can be identified:

- economic diversification. Reducing excessive dependence on the oil
sector requires the development of non-resource industries such as
manufacturing, agriculture, tourism, and information technologies.
Diversification policies should focus on attracting private investment
and supporting industrial development programs;

- institutional reforms. Strengthening institutional capacity is
essential for the effective management of natural resource revenues.
Increasing transparency, improving public governance, and strengthening
anti-corruption mechanisms will enhance investor confidence and support
sustainable development;

- human capital development. Investments in education and healthcare are
necessary to increase labor productivity and support the growth of
high-value sectors. Improving vocational and technical education will
help align workforce skills with the needs of a diversified economy;

- green economy transition. Expanding renewable energy and improving
energy efficiency can reduce dependence on fossil fuels and support
Kazakhstan's environmental sustainability goals. Such measures will also
help mitigate risks related to global carbon regulation;

- trade diversification and international cooperation. Expanding trade
partnerships and strengthening regional economic cooperation can reduce
reliance on energy exports and broaden Kazakhstan's export structure;

- digital transformation. Developing digital infrastructure and
supporting technological innovation can enhance productivity, stimulate
new industries, and strengthen the country's position in the global
economy.

Overall, these policy priorities provide a strategic framework for
reducing resource dependence and supporting long-term sustainable
development in Kazakhstan.

{\bfseries Conclusions.} This study examined the structural vulnerabilities
of Kazakhstan's resource-dependent economic model and identified key
transformation needs using a combined SWOT and PEST analytical
framework. The findings confirm the main assumptions of the natural
resource curse literature, showing that excessive reliance on natural
resources may hinder long-term development when economic diversification
and institutional capacity remain limited.

The results demonstrate that Kazakhstan's economic growth remains
strongly dependent on oil and gas revenues, which increases
vulnerability to external shocks and limits the development of other
sectors. Institutional weaknesses and insufficient diversification
further constrain the effective use of resource wealth. At the same
time, opportunities related to renewable energy, technological
development, and digital transformation create potential pathways for
economic restructuring.

From a policy perspective, the findings suggest several strategic
priorities. First, economic diversification should be accelerated by
supporting manufacturing, agriculture, tourism, and technology sectors.
Second, institutional reforms aimed at improving transparency,
governance quality, and the management of natural resource revenues are
essential. Third, investments in human capital, innovation, and green
technologies are necessary to support long-term sustainable development.

This study also highlights several limitations. The SWOT and PEST
approaches rely on qualitative assessment and therefore may involve a
certain degree of subjectivity. In addition, the analysis focuses on a
single country case, which limits the generalizability of the findings.
Future studies may expand the research through quantitative modeling and
comparative analysis of multiple resource-rich economies.

Despite these limitations, the study contributes to the literature by
demonstrating that the negative effects of resource dependence are
strongly linked to institutional quality and economic diversification.
The findings support the argument that natural resource abundance alone
does not determine economic outcomes; rather, effective governance and
appropriate economic policies play a decisive role.

Overall, Kazakhstan's experience illustrates that resource wealth can
either become a development opportunity or a structural constraint. By
strengthening institutions, diversifying its economic structure, and
investing in innovation and sustainable energy, Kazakhstan can transform
its natural resource wealth into a long-term development advantage.
\end{multicols}

\begin{center}
{\bfseries Литература}
\end{center}

\begin{refs}
1. Sachs J. D., Warner A. M. The curse of natural resources//European
Economic Review. -2001.-Vol.45(4-6).-P.827-838.
\href{https://doi.org/10.1016/S0014-2921(01)00125-8}{DOI
10.1016/S0014-2921(01)00125-8}.

2. Auty R.M. Natural resources, capital accumulation and the resource
curse // Ecological Economics.-2007.-Vol.61(4).-P.627-634. DOI
10.1016/j.ecolecon.2006.09.006.

3. Kelesbayev, D., Myrzabekkyzy, K., Bolganbayev, A., Baimaganbetov, S.
The Impact of Oil Prices on the Stock Market and Real Exchange Rate: The
Case of Kazakhstan // International Journal of Energy Economics and
Policy.-2022.-Vol.12(1).-P.163-168.DOI 10.32479/ijeep.11880.

4. Kelesbayev D., Myrzabekkyzy K., Bolganbayev A., Baimaganbetov S. The
Effects of the Oil Price Shock on Inflation: The Case of Kazakhstan //
International Journal of Energy Economics and
Policy.-2022.-Vol.12(3).P.477-481.DOI 10.32479/ijeep.13061.

5. Baimaganbetov S., Kelesbayev D., Yermankulova R., Izzatullaeva B.,
Almukhambetova B. Effects of Oil Price Changes on Regional Real Income
Per Capita in Kazakhstan: Panel Data Analysis//International Journal of
Energy Economics and Policy. -2019.- Vol.9(4). -P.356-362. DOI
10.32479/ijeep.8081.

6. Atkinson G., Hamilton K. Savings Growth and the Resource Curse
Hypothesis // World Development. -2003.-Vol.31(11). -P.1793-1807. DOI
10.1016/J.WORLDDEV.2003.05.001.

7. Baimaganbetov S., Kelesbayev D., Baibosynova G., Yermankulova R.,
Dandayeva B. The Impact of Oil Prices on the Food Inflation in
Kazakhstan // International Journal of Energy Economics and
Policy.-2021.-Vol.11(3).P.73-79. DOI 10.32479/ijeep.10844.

8. Hodler R. The Curse of Natural Resources in Fractionalized
Countries//European Economic Review.-2006.-Vol.50(6).-P.1367-1386. DOI
10.1016/j.euroecorev.2005.05.004.

9. Alexeev M., Conrad R. The Elusive Curse of Oil // The Review of
Economics and Statistics. -2009.-Vol.91(3). -P.586-598. DOI
10.1162/rest.91.3.586.

10. Brunnschweiler C. N. Cursing the Blessing? Natural Resource
Abundance,Institutions, and Economic Growth//World
Development.-2008.-Vol.36(3).-P.399-419. DOI 10.1016/j.worlddev.2007.03.004.

11. Alssadek M., Benhin J. Natural resource curse: A literature survey
and comparative assessment of regional groupings of oil-rich countries
// Resources Policy. -2023. -Vol.84: 103741. DOI
10.1016/j.resourpol.2023.103741.

12. Wang Z., Li H., Zhang Y. Sustainability education and resource curse
control in the selected resource-rich economies // Resources Policy.
-2024. -Vol.97:105274. DOI 10.1016/j.resourpol.2024.105274.

13. Herb M. No Representation without Taxation? Rents Development and
Democracy // Comparative Politics.-2005.-Vol.37(3).- P.297-317. DOI
10.2307/20072891.

14. International Monetary Fund (IMF). IMF Executive Board Concludes the
2024 Article IV Consultation with the Republic of Kazakhstan. - 2024. - URL:
\href{https://www.imf.org/en/news/articles/2025/01/30/pr25021-kazakhstan-executive-board-concludes-2024-article-iv-consult}{https://www.imf.org}. - Date of access: 10.10.2025.

15. Қазақстан Республикасы Стратегиялық жоспарлау және реформалар
агенттігінің Ұлттық статистика бюросы. Ұлттық шоттар статистикасы. -
2024. - URL: https://stat.gov.kz/industries/economy/national-accounts/.
- Date of access: 10.10.2025.
\end{refs}

\begin{center}
{\bfseries References}
\end{center}

\begin{refs}
1. Sachs J. D., Warner A. M. The curse of natural resources//European
Economic Review. -2001.-Vol.45(4-6).-P.827-838.
\href{https://doi.org/10.1016/S0014-2921(01)00125-8}{DOI
10.1016/S0014-2921(01)00125-8}.

2. Auty R.M. Natural resources, capital accumulation and the resource
curse // Ecological Economics.-2007.-Vol.61(4).-P.627-634. DOI
10.1016/j.ecolecon.2006.09.006.

3. Kelesbayev, D., Myrzabekkyzy, K., Bolganbayev, A., Baimaganbetov, S.
The Impact of Oil Prices on the Stock Market and Real Exchange Rate: The
Case of Kazakhstan // International Journal of Energy Economics and
Policy.-2022.-Vol.12(1).-P.163-168.DOI 10.32479/ijeep.11880.

4. Kelesbayev D., Myrzabekkyzy K., Bolganbayev A., Baimaganbetov S. The
Effects of the Oil Price Shock on Inflation: The Case of Kazakhstan //
International Journal of Energy Economics and
Policy.-2022.-Vol.12(3).P.477-481.DOI 10.32479/ijeep.13061.

5. Baimaganbetov S., Kelesbayev D., Yermankulova R., Izzatullaeva B.,
Almukhambetova B. Effects of Oil Price Changes on Regional Real Income
Per Capita in Kazakhstan: Panel Data Analysis//International Journal of
Energy Economics and Policy. -2019.- Vol.9(4). -P.356-362. DOI
10.32479/ijeep.8081.

6. Atkinson G., Hamilton K. Savings Growth and the Resource Curse
Hypothesis // World Development. -2003.-Vol.31(11). -P.1793-1807. DOI
10.1016/J.WORLDDEV.2003.05.001.

7. Baimaganbetov S., Kelesbayev D., Baibosynova G., Yermankulova R.,
Dandayeva B. The Impact of Oil Prices on the Food Inflation in
Kazakhstan // International Journal of Energy Economics and
Policy.-2021.-Vol.11(3).P.73-79. DOI 10.32479/ijeep.10844.

8. Hodler R. The Curse of Natural Resources in Fractionalized
Countries//European Economic Review.-2006.-Vol.50(6).-P.1367-1386. DOI
10.1016/j.euroecorev.2005.05.004.

9. Alexeev M., Conrad R. The Elusive Curse of Oil // The Review of
Economics and Statistics. -2009.-Vol.91(3). -P.586-598. DOI
10.1162/rest.91.3.586.

10. Brunnschweiler C. N. Cursing the Blessing? Natural Resource
Abundance,Institutions, and Economic Growth//World
Development.-2008.-Vol.36(3).-P.399-419. DOI 10.1016/j.worlddev.2007.03.004.

11. Alssadek M., Benhin J. Natural resource curse: A literature survey
and comparative assessment of regional groupings of oil-rich countries
// Resources Policy. -2023. -Vol.84: 103741. DOI
10.1016/j.resourpol.2023.103741.

12. Wang Z., Li H., Zhang Y. Sustainability education and resource curse
control in the selected resource-rich economies // Resources Policy.
-2024. -Vol.97:105274. DOI 10.1016/j.resourpol.2024.105274.

13. Herb M. No Representation without Taxation? Rents Development and
Democracy // Comparative Politics.-2005.-Vol.37(3).- P.297-317. DOI
10.2307/20072891.

14. International Monetary Fund (IMF). IMF Executive Board Concludes the
2024 Article IV Consultation with the Republic of Kazakhstan. - 2024. - URL:
\href{https://www.imf.org/en/news/articles/2025/01/30/pr25021-kazakhstan-executive-board-concludes-2024-article-iv-consult}{https://www.imf.org}. - Date of access: 10.10.2025.

15. Қazaқstan Respublikasy Strategijalyқ zhosparlau zhәne reformalar
agenttіgіnің Ұlttyқ statistika bjurosy. Ұlttyқ shottar statistikasy. -
2024. - URL: https://stat.gov.kz/industries/economy/national-accounts/.
- Date of access: 10.10.2025. {[}in Kazakh{]}.
\end{refs}

\begin{info}
\hspace{1em}\emph{{\bfseries Information about the authors}}

Koptayeva G.P. - Candidate of Economic Sciences, Senior Lecturer,
Miras University, Shymkent, Kazakhstan, e-mail: asel\_4747@mail.ru;

Mukhamedkhanova A.B. - Doctoral Student, M. Auezov South Kazakhstan
University, Shymkent, Kazakhstan, e-mail: Dia-2808@mail.ru;

Issakhmetova A.N. - Candidate of Economic Sciences, Associate
Professor, Central Asian Innovative University, Shymkent, Kazakhstan,
e-mail: i\_a\_n@inbox.ru;

Myrzalieva G.A - Candidate of Economic Sciences, Senior Lecturer,
Miras University, Shymkent, Kazakhstan, e-mail:
k.i.myrzalieva@yandex.ru;

Mutaliyeva A.A - PhD, Senior Lecturer, Regional Innovation University,
Shymkent, Kazakhstan, e-mail: Alua012@mail.ru;

Kelesbayev D.N. - PhD, Professor, Ahmed Yasawi University, Turkestan,
Kazakhstan, e-mail: dinmukhamed.kelesbayev@ayu.edu.kz.

\hspace{1em}\emph{{\bfseries Сведения об авторах}}

Коптаева Г.П. - кандидат экономических наук, старший преподаватель,
Университет Мирас, Шымкент, Казахстан, е-mail: asel\_4747@mail.ru;

Мухамедханова А.Б. - докторант, Южно-Казахстанский университет им.
М.Ауэзова, Шымкент, Казахстан, е-mail: Dia-2808@mail.ru;

Исахметова А.Н. - кандидат экономических наук, доцент,
Центрально-Азиатский инновационный университет, Шымкент, Казахстан,
е-mail: i\_a\_n@inbox.ru;

Мырзалиева Г.А. - кандидат экономических наук, старший преподаватель,
Университет Мирас, Шымкент, Казахстан, е-mail:
k.i.myrzalieva@yandex.ru;

Муталиева А.А. - PhD, старший преподаватель, Региональный
инновационный университет, Шымкент, Казахстан, е-mail:
Alua012@mail.ru;

Келесбаев Д.Н. - PhD, профессор, Университет Ахмеда Ясави, Туркестан,
Казахстан, е-mail: dinmukhamed.kelesbayev@ayu.edu.kz.
\end{info}
